\subsubsection*{Solution}
\paragraph*{Step-by-Step Derivation}

\subparagraph*{1. Magic-wavelength condition}

For a monochromatic trap with intensity $I$, the AC Stark shift of an atomic state $v$ is

\[
\Delta E_v(\omega) = -\frac{\alpha_v(\omega)I}{2\epsilon_0c},
\]

where $\alpha_v(\omega)$ is its dynamic polarizability. The light-induced shift of an imaging transition is therefore

\[
\delta\nu(\omega) = -\frac{I}{2h\epsilon_0c} \left[\alpha_e(\omega)-\alpha_g(\omega)\right].
\]

Thus a magic wavelength satisfies

\[
{\alpha_e(\omega_m)=\alpha_g(\omega_m)}.
\]

The $^1S_0$ ground state has $J=0$ and only a scalar polarizability. The $^3P_1$ state has both scalar and tensor contributions:

\[
\alpha_{Jm}(\omega,\theta) = \alpha^{(0)}(\omega) + \alpha^{(2)}(\omega) \frac{3m^2-J(J+1)}{J(2J-1)} \frac{3\cos^2\theta-1}{2},
\]

where $\theta$ is the angle between the trap polarization and the quantization axis. Consequently, the numerical magic wavelengths are polarization-geometry dependent. Below, the standard linearly polarized imaging geometries used experimentally are assumed.

{\textit{Comment: The wavelengths value quoted below are taken from different papers with different experimental configurations, so the geometry should be stated separately for each reported wavelength rather than treated as one common ``assumed'' geometry.}}

Define

\[
\Delta\alpha_0(\omega) = \alpha_{^3P_1,m_J=0}(\omega)-\alpha_g(\omega),
\]

\[
\Delta\alpha_1(\omega) = \alpha_{^3P_1,|m_J|=1}(\omega)-\alpha_g(\omega).
\]

\subparagraph*{2. $^{174}\mathrm{Yb}$}

Since $^{174}\mathrm{Yb}$ has nuclear spin $I=0$, there is no hyperfine structure.

For the $\pi$ branch,

\[
|^1S_0,J=0,m_J=0\rangle \longrightarrow |^3P_1,J'=1,m'_J=0\rangle,
\]

so the magic condition is

\[
\Delta\alpha_0(\lambda_m)=0.
\]

The measured crossing is near

\[
{\lambda_m\simeq532\ \mathrm{nm}},
\]

the wavelength employed for narrow-line cooling and imaging of $^{174}\mathrm{Yb}$. The residual measured differential shift at nominal $532$ nm is about $1.6\%$ of the trap depth, so "532 nm" denotes the experimentally useful near-magic condition rather than a sub-nanometre determination. \href{\detokenize{https://arxiv.org/abs/1810.10517}}{Saskin et al.}

{\textit{Comment: The value is experimentally reported near-magic taken from the literature}}
For the $\sigma^\pm$ branch,

\[
m_J=0\longrightarrow m'_J=\pm1,
\]

and the condition becomes

\[
\Delta\alpha_1(\lambda_m)=0.
\]

The same electronic polarizability is sampled by the stretched $^{171}\mathrm{Yb}$ transition discussed below, giving a second crossing near

\[
{\lambda_m\simeq483\ \mathrm{nm}}.
\]

This value is currently reported only as "near 483 nm," so additional digits are not justified.

{\textit{Comment: The 483nm value is for the 171Yb transition, and the 174Yb entry is INFERRED from the shared electronic polarizability. }}

\subparagraph*{3. $^{171}\mathrm{Yb}$ hyperfine structure}

For $^{171}\mathrm{Yb}$, $I=1/2$. The relevant imaging manifold is

\[
^1S_0:\ F=\frac12, \qquad ^3P_1:\ F'=\frac32.
\]

The stretched excited state is

\[
\left|\frac32,\frac32\right\rangle = |m_J=1,m_I=\tfrac12\rangle.
\]

It therefore has the same electronic polarizability as the $|m_J|=1$ state of $^{174}\mathrm{Yb}$. Consequently,

\[
\Delta\alpha_{F'=3/2,|m_F'|=3/2} = \Delta\alpha_1.
\]

The corresponding cycling transition is

\[
\left|\frac12,\pm\frac12\right\rangle \longrightarrow \left|\frac32,\pm\frac32\right\rangle,
\]

with $\Delta m_F=\pm1$, hence it is a $\sigma^\pm$ transition. Experimentally its magic wavelength is near $483$ nm, with the trap polarization perpendicular to the magnetic field. \href{\detokenize{https://arxiv.org/abs/2305.19119}}{Norcia et al.}

For the nonstretched state, Clebsch-Gordan decomposition gives

\[
\left|\frac32,\frac12\right\rangle = \sqrt{\frac23} |m_J=0,m_I=\tfrac12\rangle + \sqrt{\frac13} |m_J=1,m_I=-\tfrac12\rangle.
\]

Since the electric-dipole trapping interaction does not act on the nuclear spin,

\[
\alpha_{F'=3/2,|m_F'|=1/2} = \frac23\alpha_0+\frac13\alpha_1.
\]

The magic condition is therefore

\[
\frac23\alpha_0+\frac13\alpha_1=\alpha_g.
\]

Subtracting $\alpha_g$ yields

\[
\frac23\Delta\alpha_0+\frac13\Delta\alpha_1=0,
\]

or equivalently,

\[
{\frac{\Delta\alpha_0}{\Delta\alpha_1}=-\frac12}.
\]

The measured crossing of this ratio occurs at

\[
{\lambda_m=486.78\pm0.10\ \mathrm{nm}}
\]

in vacuum. At $486.78$ nm the measured differential shift was only $0.0001(6)$ times the ground-state trap depth. For the standard same-$m_F$ imaging branch,

\[
\left|\frac12,\pm\frac12\right\rangle \longrightarrow \left|\frac32,\pm\frac12\right\rangle,
\]

$\Delta m_F=0$, so this is classified as a $\pi$ transition. \href{\detokenize{https://arxiv.org/abs/2112.06799}}{Ma et al.}


The $532$ nm condition is not magic for $^{171}\mathrm{Yb}$ because its $|m_F'|=1/2$ state contains both $m_J=0$ and $|m_J|=1$ components.
{\textit{Comment: These four values might not be the only 4 possible set of crossings in the range, for example: "Magic wavelengths for the 6s ² ¹ S 0 - 6s6p ³ P 1 transition in ytterbium atom"}}